\section{Build It: A Quantum-Inspired GNN Layer}
\label{sec:buildit}

We now assemble a concrete, param-matched model. The design principle is that all
three operators share architecture and differ \emph{only} in $\Prop$, so any
performance gap is attributable to the operator and not to capacity.

\subsection{The shared layer}
Given layer embeddings $\Feat^{(\ell)}$, every model computes
\begin{equation}
\Feat^{(\ell+1)}=\sigma\!\big(\Prop^{(\ell)}\,\Feat^{(\ell)}\,\Wt^{(\ell)}+\mathbf{b}^{(\ell)}\big),
\end{equation}
and differs only in the operator:
\begin{align}
\text{GCN:}&\quad \Prop^{(\ell)}=\hat{\Adj}\ \text{(fixed)},\\
\text{Heat:}&\quad \Prop^{(\ell)}=\Eig\,\mathrm{diag}(e^{-t_\ell\lambda})\,\Eig^{\top},\\
\text{QI:}&\quad \Prop^{(\ell)}=\big|\Eig\,\mathrm{diag}(e^{-it_\ell\lambda})\,\Eig^{*}\big|^2 .
\end{align}
The heat and quantum-inspired (QI) operators add one learnable scalar $t_\ell$ per
layer, which is negligible against the weight matrices, so the three models have
essentially identical parameter counts. Because $\Prop$ is built from the eigenvalues
$\lambda$ with differentiable functions, $t_\ell$ trains by ordinary backpropagation;
a single eigendecomposition of $\Lap$ per graph is reused across layers and operators.

\subsection{Algorithm}
Algorithm~\ref{alg:qignn} gives the QI-GNN forward pass; the heat and GCN variants
are the same template with their operator substituted. Fig.~\ref{fig:qignn} shows the
data flow.

\begin{algorithm}[t]
\caption{QI-GNN forward pass (one graph)}
\label{alg:qignn}
\begin{algorithmic}[1]
\REQUIRE features $\mathbf{X}$; eigenpairs $(\lambda,\Eig)$ of $\Lap$;
parameters $\{t_\ell,\Wt^{(\ell)},\mathbf{b}^{(\ell)}\}_{\ell=0}^{K-1}$
\STATE $\Feat\leftarrow\mathbf{X}$
\FOR{$\ell=0$ to $K-1$}
  \STATE $\Prop\leftarrow\big|\Eig\,\mathrm{diag}(e^{-i\,t_\ell\lambda})\,\Eig^{*}\big|^2$
         \hfill\COMMENT{quantum-walk probability kernel}
  \STATE $\Feat\leftarrow\sigma\big(\Prop\,\Feat\,\Wt^{(\ell)}+\mathbf{b}^{(\ell)}\big)$
\ENDFOR
\RETURN $\mathrm{readout}(\Feat)$
\end{algorithmic}
\end{algorithm}

\begin{figure}[t]
\centering
\resizebox{\columnwidth}{!}{%
\begin{tikzpicture}[node distance=5mm]
% main horizontal data flow (clean line, no box in the way)
\node[box] (x) {features\\$\mathbf{X}$};
\node[hot, right=12mm of x] (prop) {$\Prop=|\Eig e^{-it_\ell\lambda}\Eig^{*}|^2$};
\node[box, right=12mm of prop] (mix) {$\sigma(\Prop\Feat\Wt_\ell)$};
\node[plain, right=10mm of mix] (out) {readout\\$\hat{\mathbf{y}}$};
% eigendecomposition feeds the operator from below
\node[plain, below=8mm of prop] (eig) {eig of $\Lap$: $(\lambda,\Eig)$};

\draw[flow] (x) -- (prop);
\draw[flow] (eig) -- (prop);
\draw[flow] (prop) -- (mix);
\draw[flow] (mix) -- (out);
% recurrence over K layers: arc well above the boxes, label above the arc
\draw[flow, dashed] (mix.north) .. controls +(0,10mm) and +(0,10mm) .. (prop.north);
\node[font=\scriptsize, text=ink, above=1pt]
  at ($(prop.north)!0.5!(mix.north) + (0,10mm)$)
  {$\times K$ layers, learnable $t_\ell$};
\end{tikzpicture}}
\caption{QI-GNN data flow (Algorithm~\ref{alg:qignn}). The quantum-walk kernel is
built once from the Laplacian spectrum (fed in from below) and reused across $K$ layers;
the heat and GCN variants substitute their operator into the same template, keeping the
models param-matched.}
\label{fig:qignn}
\end{figure}


\subsection{Design variants}
The quantum-walk family has design knobs worth knowing, which we probe empirically in
Section~\ref{sec:casestudy}: the walk can run on the adjacency $e^{-it\Adj}$ or the
normalized Laplacian $e^{-it\Lap^{\mathrm{sym}}}$ instead of $\Lap$, and a layer can
concatenate several learnable times $\{t_{\ell,m}\}$ for a multi-time readout. These
are alternatives a practitioner would naturally try when seeking an advantage, so we
test them rather than assume one of them is best.

\subsection{A worked micro-example}
To make the operator concrete, take the $3$-node path $1\!-\!2\!-\!3$ with combinatorial
Laplacian
\begin{equation}
\Lap=\begin{pmatrix} 1 & -1 & 0\\ -1 & 2 & -1\\ 0 & -1 & 1\end{pmatrix},
\quad \text{eigenvalues } \lambda\in\{0,1,3\}.
\end{equation}
Seeding the centre node and propagating for a small time $t$, the heat kernel keeps the
mass concentrated near the centre and merely leaks symmetrically outward, whereas the
quantum-walk probability $|e^{-it\Lap}|^2$ develops oscillatory amplitude that, as $t$
grows, pushes weight onto the two endpoints faster than diffusion would. On this tiny
graph the difference is modest, but it is the same mechanism that, scaled to a constellation
of large diameter, separates a global operator from a local one and produces the far-node
behavior measured in Section~\ref{sec:casestudy}. In code, the entire layer is the three
lines of Algorithm~\ref{alg:qignn}: form the eigenpairs once, build $\Prop$ from a
learnable $t_\ell$, and apply $\sigma(\Prop\Feat\Wt_\ell)$.

\subsection{Cost}
The three operators differ in cost as well as physics. The GCN operator $\hat{\Adj}$ is
sparse, so a layer costs $O(|E_t|\,d)$ for $d$-dimensional features. The heat and
quantum-walk kernels are dense $N\times N$ matrices built from one symmetric
eigendecomposition per graph, an $O(N^3)$ step that dominates at our scale but is shared
across layers and across all spectral operators, so it is amortized. For the hundreds of
nodes in the case study this is a fraction of a second on a CPU. It does not scale to
thousands of satellites, which is why Section~\ref{sec:roadmap} turns to
Krylov-subspace and Chebyshev approximations that compute the action $e^{-it\Lap}\Feat$
without ever forming the full kernel or its spectrum. The practical takeaway is that the
quantum-inspired operator carries no training-time penalty relative to classical
diffusion: they share the same bottleneck.

\subsection{A practical recipe}
To apply this framework to a new dynamic-wireless task, the following recipe captures the
whole workflow. (1) \emph{Build the graph}: define $G_t$ from the network state, the
adjacency by physical or logical links, and choose features that the operator can
propagate (for a potential field, a one-hot source plus a normalized degree suffices).
(2) \emph{Pick the operator by range}: if the task is short-range, start with a $K$-layer
GCN; if it is long-range on a large-diameter graph, use a global operator and treat the
classical heat kernel as the default. (3) \emph{Add the quantum-inspired option as an
ablation}, not as the hero: include $|e^{-it\Lap}|^2$ and, if you pursue it, sweep its
base matrix (adjacency, Laplacian, normalized Laplacian) as in Exp~B2. (4) \emph{Keep it
param-matched}: vary only the operator, holding depth, width, and weight shapes fixed.
(5) \emph{Evaluate honestly}: multiple seeds with mean and standard deviation, a
scale-free metric (AURC), honest denominators, and a verdict that treats any within-noise
ranking as a tie. (6) \emph{Scale with care}: replace the $O(N^3)$ eigendecomposition by a
Chebyshev or Krylov approximation of $e^{-it\Lap}\Feat$ before going to thousands of
nodes. The case study in Section~\ref{sec:casestudy} is a worked instance of exactly this
recipe.
